% ---------- page and fonts ----------
\usepackage[a4paper,margin=2.6cm]{geometry}
\usepackage[T1]{fontenc}
\usepackage{lmodern}
\usepackage{microtype}
\usepackage[dvipsnames]{xcolor}
\usepackage{enumitem}
\usepackage{booktabs}
\setlength{\parindent}{0pt}
\setlength{\parskip}{0.55em}

% ---------- mathematics ----------
\usepackage{amsmath,amssymb,amsthm,mathtools,bm}

% Shared notation. A book's own symbols go in notation.tex (loaded at the end if present).
\DeclareMathOperator{\E}{E}                     % write \E_{y\mid\vect{x}}\!\left[\cdot\right]
\DeclareMathOperator{\Var}{Var}
\DeclareMathOperator{\Cov}{Cov}
\DeclareMathOperator{\Bias}{Bias}
\DeclareMathOperator{\MSE}{MSE}
\DeclareMathOperator{\tr}{tr}
\DeclareMathOperator{\rank}{rank}
\DeclareMathOperator*{\argmax}{arg\,max}
\DeclareMathOperator*{\argmin}{arg\,min}
\newcommand{\given}{\mid}
\newcommand{\iid}{\overset{\text{iid}}{\sim}}
\newcommand{\R}{\mathbb{R}}
\newcommand{\ind}[1]{\mathbf{1}\{#1\}}          % indicator
\newcommand{\T}{^{\mathsf{T}}}                  % transpose

% ---------- boxes and theorem-like environments ----------
\usepackage{tcolorbox}
\tcbuselibrary{skins,breakable}
\definecolor{defc}{HTML}{1F4F8F}
\definecolor{thmc}{HTML}{1F7A4D}
\definecolor{intc}{HTML}{7A5A12}
\definecolor{prec}{HTML}{5B6068}
\definecolor{exc}{HTML}{6A4C93}
\definecolor{trapc}{HTML}{A23B2A}

% ---------- links (before the theorem types, so cleveref names each type correctly) ----------
\usepackage[hidelinks]{hyperref}   % PDF outline only; the text has no links (PLAN.md, D6)
\usepackage[capitalise,noabbrev]{cleveref}
\crefname{result}{Result}{Results}
\crefname{steps}{Step}{Steps}

% Body starts on a new line under the head (so a list never sits beside the title); upright text throughout.
\newtheoremstyle{slmthm}{0pt}{0pt}{\normalfont}{}{\bfseries}{.}{\newline}{}
\newtheoremstyle{slmrem}{0pt}{0pt}{\normalfont}{}{\itshape}{.}{\newline}{}
\usepackage{thmtools}   % declares the types so cleveref names each one correctly despite the shared counter
\declaretheorem[style=slmthm,numberwithin=chapter]{theorem}
\declaretheorem[style=slmthm,sibling=theorem]{result}
\declaretheorem[style=slmthm,sibling=theorem]{lemma}
\declaretheorem[style=slmthm,sibling=theorem]{corollary}
\declaretheorem[style=slmthm,sibling=theorem]{definition}
\declaretheorem[style=slmthm,sibling=theorem]{example}
\declaretheorem[style=slmrem,sibling=theorem]{remark}

% amsthm's proof, with "Proof." on its own line
\makeatletter
\renewenvironment{proof}[1][\proofname]{\par\pushQED{\qed}\normalfont\topsep6\p@\@plus6\p@\relax\trivlist\item[\hskip\labelsep\itshape#1\@addpunct{.}]\mbox{}\par\nobreak}{\popQED\endtrivlist\@endpefalse}
\makeatother

\tcbset{slm/.style={enhanced,before upper={\setlength{\parskip}{0.55em}},breakable,boxrule=0.5pt,arc=1.5pt,left=6pt,right=6pt,top=4pt,bottom=4pt,before skip=10pt,after skip=10pt}}
\tcolorboxenvironment{definition}{slm,colback=defc!4,colframe=defc!70}
\tcolorboxenvironment{theorem}{slm,colback=thmc!5,colframe=thmc!75}
\tcolorboxenvironment{result}{slm,colback=thmc!5,colframe=thmc!75}
\tcolorboxenvironment{lemma}{slm,colback=thmc!5,colframe=thmc!75}
\tcolorboxenvironment{corollary}{slm,colback=thmc!5,colframe=thmc!75}
\tcolorboxenvironment{example}{slm,colback=exc!4,colframe=exc!60}
% A box never starts in the last few lines of a page (its title would sit alone there): the needspace trick, skipped right after a heading.
\makeatletter\newcommand{\keepwithnext}[1]{\par\if@nobreak\else\begingroup\setlength{\dimen@}{#1}\vskip\z@\@plus\dimen@\penalty-100\vskip\z@\@plus-\dimen@\vskip\dimen@\penalty9999\vskip-\dimen@\vskip\z@skip\endgroup\fi}\makeatother
\AddToHook{env/definition/before}{\keepwithnext{7\baselineskip}}
\AddToHook{env/theorem/before}{\keepwithnext{7\baselineskip}}
\AddToHook{env/result/before}{\keepwithnext{7\baselineskip}}
\AddToHook{env/lemma/before}{\keepwithnext{7\baselineskip}}
\AddToHook{env/corollary/before}{\keepwithnext{7\baselineskip}}
\AddToHook{env/example/before}{\keepwithnext{7\baselineskip}}
% a heading needs room for itself and the start of what follows (a breakable box after it can otherwise strand it)
\AddToHook{cmd/section/before}{\keepwithnext{10\baselineskip}}
\AddToHook{cmd/subsection/before}{\keepwithnext{8\baselineskip}}

\newtcolorbox{question}{slm,colback=exc!3,colframe=exc!50,title={The question},fonttitle=\bfseries\small}
\newtcolorbox{setup}[1][Setup]{slm,colback=defc!2,colframe=defc!45,title={#1},fonttitle=\bfseries\small}
\newtcolorbox{trap}{slm,colback=trapc!4,colframe=trapc!70,title={Trap},fonttitle=\bfseries\small}

% Two-column derivation (PLAN.md, D5): maths on the left, the reason on the right.
%   \begin{deriv} \dl{maths}{reason} ... \end{deriv}
\newenvironment{deriv}{\par\nopagebreak\smallskip\noindent\renewcommand{\arraystretch}{1.3}%
  \begin{tabular}{@{}p{0.58\linewidth}@{\hspace{0.03\linewidth}}p{0.39\linewidth}@{}}}{\end{tabular}\par\smallskip}
\newcommand{\dl}[2]{\mbox{$\displaystyle #1$} & \leavevmode\raggedright\small\color{thmc!80!black}#2\tabularnewline[1.2ex]}

% Numbered derivation steps: Step-1, Step-2, ...
\newlist{steps}{enumerate}{1}
\setlist[steps]{label=\textbf{Step-\arabic*.},ref=\arabic*,leftmargin=*,align=left,itemsep=0.4em}
\setlist[itemize]{itemsep=0.15em,topsep=0.3em}

% ---------- small markers ----------
% first, defining use of a term: bold + a glossary entry (order checked by tools/check.py)
\usepackage{makeidx}\makeindex
\renewcommand{\indexname}{Glossary}
\newcommand{\dfn}[1]{\textbf{#1}\index{#1}}

\InputIfFileExists{notation}{}{}

